\section{Related Work}
\label{sec:related}

\subsection{Knowledge Tracing Models}
Bayesian Knowledge Tracing (BKT) models mastery of each skill with a
two-state hidden Markov model~\cite{corbett1994bkt}. Deep Knowledge
Tracing (DKT) replaced this with a recurrent neural network, improving
predictive accuracy at the cost of per-skill interpretability~%
\cite{piech2015dkt}. Attention-based models followed: SAKT applies
self-attention over the interaction sequence~\cite{pandey2019sakt}, and
AKT introduces a monotonic attention mechanism with a
Rasch-model-based embedding~\cite{ghosh2020akt}. More recently, simpleKT
was proposed as a strong yet simple baseline that is difficult to beat
across domains~\cite{liu2023simplekt}. The pyKT toolkit provides
standardised preprocessing and reference implementations of these models
across several public datasets, enabling transparent
comparison~\cite{liu2022pykt}.
% \note{These four (DKT, SAKT, AKT, simpleKT) are your RO1 baselines.}

\subsection{Programming Knowledge Tracing}
A growing line of work adapts KT to programming, where each interaction
carries source code rather than only a correctness label. Code-DKT
extends DKT with an attention mechanism that extracts domain-specific
code features~\cite{shi2022codedkt}. Other work interacts programming
skills with student code representations such as code2vec and
ASTNN~\cite{pktattn2021}. Test-case-informed KT (TIKTOC) incorporates
fine-grained test-case outcomes for open-ended coding
tasks~\cite{tiktoc2025}. KCGen-KT automatically generates knowledge
components from solution ASTs and traces them~\cite{kcgenkt2026}.
HELP-DKT models introductory Python programming with an interpretable
cognitive model built on deep knowledge tracing~\cite{helpdkt2022}. Our proposed variant (Section~%
\ref{sec:methodology}) builds on this direction by combining
knowledge-component tags with code-structure features on top of a strong
attention baseline.

\subsection{KT as a Simulator for Instructional Sequencing}
Beyond prediction, KT models are increasingly used as simulators of
learners to train or evaluate instructional-sequencing
policies~\cite{piech2015dkt}. Reinforcement-learning approaches optimise
sequencing against such simulators, and recent work (e.g., multi-step
simulator training) aims to make these simulators more robust. However,
the coupling between a model's predictive accuracy and its usefulness as
a sequencing simulator is known to be weak in general KT and largely
unexamined for programming---the gap this work targets.
% \todo{Add AdvKT / RL-DKT / deep-RL sequencing citations to references.bib and cite here.}

\subsection{Public Datasets}
For programming, the CodeWorkout dataset from the CSEDM 2021 Data
Challenge provides Java CS1 submissions to 50 problems with
knowledge-component tags~\cite{csedm2021}. FalconCode provides a
multi-year corpus of introductory Python submissions with per-problem
skills and test cases~\cite{falconcode2023}. For the mathematics
contrast we use ASSISTments / Junyi as provided by
pyKT~\cite{liu2022pykt}.
